\documentclass[10pt,a4paper]{amsart}
\usepackage[margin=19mm]{geometry}
\usepackage{amsmath,amssymb,mathtools}
\usepackage[scr=rsfs,cal=euler]{mathalfa}
\usepackage{xfrac}
\usepackage[unicode,hidelinks,bookmarks=false]{hyperref}
\newcommand{\ie}{i.e.\ }
\newcommand{\cf}{cf.\ }
\newcommand{\resp}{respectively\ }
\newcommand{\Subsection}{\subsection}
\providecommand{\nobreakdash}{\nobreak\hbox{-}}
\setlength{\emergencystretch}{2em}
\multlinegap0pt
\usepackage{bm}
\usepackage{bbm}
\usepackage{dsfont}
\usepackage{xspace}
\usepackage{cancel}
\usepackage{mathtools}
\usepackage{amsthm}
\makeatletter
\@ifundefined{theorem}{\newtheorem{theorem}{Theorem}}{}
\@ifundefined{corollary}{\newtheorem{corollary}[theorem]{Corollary}}{}
\makeatother
\def\Cscr{{\mathscr C}}  
\def\HH{{\mathbb H}}
\def\Hcal{{\mathcal H}} 
\def\Ibold{\mathbf {I}}
\def\Ical{{\mathcal I}} 
\def\Nbold{{\mathbf N}} 
\def\g{\gamma}
\def\pt{{\scriptscriptstyle\bullet}}
\def\tri{\triangle}
\newcommand\Htp{\operatorname{Htp}}
\newcommand\att{\operatorname{att}}
\newcommand\conf{\mathscr{C}\!\mathit{onf}\!}
\newcommand\im{\operatorname{Im}}
\title[The configuration functor of a punctured space]{The configuration functor of a punctured space}
\author{Eduard Looijenga, Andreas Stavrou}
\date{Source: arXiv:2507.14366v1 (CC BY 4.0)}
\begin{document}
\maketitle

\noindent\emph{Source and licence.} The configuration functor of a punctured space. Version arXiv:2507.14366v1 (CC BY 4.0), distributed under the Creative Commons Attribution 4.0 International licence, as recorded in the arXiv licence metadata for this exact version and shown on the abstract page https://arxiv.org/abs/2507.14366v1. Full rights notice: licenses/arxiv-2507.14366-CC-BY-4.0.txt.\par\smallskip

\noindent\textit{Editorial scope.} Counted unit: the Corollary whose source label is \texttt{cor:zeta-generation} in the article (source0.tex lines 927--929, source label \texttt{cor:zeta-generation}), delivered together with the article's own complete original proof of it (source0.tex lines 930--948, one \texttt{proof} environment of 19 lines). No mathematics is written, repaired or re-derived: statement and proof are the article's own text line for line. Required context retained uncounted: eqn:tensormap1 (lines 549--553); thm:generation (lines 555--559). Every definition, display and label of those lines is retained unchanged. Omitted from this package and not redistributed: 1--548, 554--554, 560--926, 949--1421. Nothing inside a retained excerpt is omitted or altered: every retained body is delivered exactly as the article's lines are written, so no omission receipt of any kind accompanies this package. Citations: the retained excerpts carry no citation command at all, so no bibliography entry of the article is transcribed and no reference is redistributed. Reference adaptation: none was needed and none was made; every reference inside the delivered unit resolves inside it, so no unresolved reference or citation is produced. Printed slips kept literal: no printed word, symbol, hypothesis or number of the counted statement or of its proof was altered. Layout and class: the article's own class is \texttt{amsart} with packages including amsmath, amsthm, amsbsy, amsfonts, amssymb, txfonts, shuffle, ulem, stmaryrd, mathrsfs, graphicx, supertabular, amscd, tikz-cd; the wrapper is the lane's pinned \texttt{amsart} editorial wrapper at 10pt on A4, which restates the theorem-like environments the retained text needs and adds the article's own macro definitions that the retained text needs (\texttt{\textbackslash Cscr}, \texttt{\textbackslash HH}, \texttt{\textbackslash Hcal}, \texttt{\textbackslash Htp}, \texttt{\textbackslash Ibold}, \texttt{\textbackslash Ical}, \texttt{\textbackslash Nbold}, \texttt{\textbackslash att}, \texttt{\textbackslash conf}, \texttt{\textbackslash g}, \texttt{\textbackslash im}, \texttt{\textbackslash pt}, \texttt{\textbackslash tri}), with extra wrapper packages (bm, bbm, dsfont, xspace, cancel, mathtools). The delivered numbering is the unit's own. Assets: no retained span contains a figure environment, a graphics-inclusion command, a PDF-page inclusion, a table or a TikZ picture; no image file is copied into this package, the package declares no asset dependency and no image file is redistributed with it. Licence: version arXiv:2507.14366v1 of this article is distributed under the Creative Commons Attribution 4.0 International licence, as recorded in the arXiv licence metadata for this exact version and shown on the abstract page https://arxiv.org/abs/2507.14366v1. A full rights notice is delivered at licenses/arxiv-2507.14366-CC-BY-4.0.txt. Characters: every retained excerpt is pure ASCII, so the delivered engine, XeLaTeX with its default Latin Modern text font, reports no missing character.
\begin{multline}\label{eqn:tensormap1}
\tri_U^{(\Nbold_1, \dots,\Nbold_r)}: \Ical_{U}|_{n_1}\otimes\Ical_{U}|_{n_2}\otimes\cdots\otimes \Ical_{U}|_{n_r}
\xrightarrow{\triangle_U^{\Nbold_1}\otimes\cdots\otimes \triangle_U^{\Nbold_r}}\\
\to \Hcal_{N_1}(U)\otimes\Hcal_{N_2}(U)\otimes\cdots \otimes\Hcal_{N_r}(U)\to \Hcal_N(U).
\end{multline}

\begin{theorem}\label{thm:generation}
Assume $n>0$. Then the images of the  maps \eqref{eqn:tensormap1} generate $\Hcal_N(U)$. We  already get a set of generators by evaluating these maps on  tensors of the form $(\g_{c_1}-1)|_{n_1}\otimes\cdots\otimes (\g_{c_r}-1)|_{n_r}$, where each $c_i$ is a $1$-cell of $U^*$ and $\g_{c_i}\in \pi_U$ is defined by an orientation of $c_i$.

In particular,  the group $\Htp(U^*,*)$ of self-homotopy equivalences acts on $\Hcal_N(U)$  through its action on $\Ical_{U}|_n$.
\end{theorem} 

\begin{corollary}\label{cor:zeta-generation}
Assume $\dim U_o=1$,  so that $U_o^*$ is the  $1$-skeleton of $U$. If  a homotopy class $f_o\in \Htp(U_o,*)$ fixes the element  $\zeta\in \pi_{U_o}$ defined by the attaching map,  then $f_o$ extends to a  $f\in \Htp(U,*)$ whose action on  $H^{cl}_k(\conf_N(U))$ is  through its action  on $\Ical_{U}|_{n-k}$.
\end{corollary}

\begin{proof}
In this case, $\Cscr_K(U_o)_\pt$ is a single term complex concentrated in degree $|K|$ and then equal to $\Hcal_K(U_o)$.

Since $f_o$ fixes $\zeta$, it can be represented by a continuous map $F_o:U_o\to U_o$ together with a homotopy between $\att$ and $F_o\circ att$. 
Perform this homotopy (relative to the base point) on the closure of the strip  $0<\im(z)\le 1$ of $\HH_+$ in $U$ and extend it by the map $z\mapsto z-\sqrt{-1}$ on $\im(z)\ge 1$ to obtain the cellular map $F: U\to U$. 
By construction, the map $F$  preserves the total order we imposed on $\HH_+$. These conditions are precisely what we need so that $F$ acts cellularly on the pair  $(U^{*N}, D_N(U))$.  It then acts functorially on  $\Cscr_N(U)_\pt$ in a way that respects the monodial structure. On the subcomplex $\Cscr_{N_0}(U)_\pt=\Hcal_{N_0}(U_o)$ the action is via $f_o$. 
The action of $F$ on the other  tensor factor $\Hcal_{N_i}[1]$ (which defines   in general not a subcomplex) is trivial: this group  coincides with 
$\Cscr_N(\HH_+)_\pt$ on which $F$ acts through the quotient $U/U_o\cong (\HH_+)^*$, but on this quotient 
$F$ is homotopic to the identity via the lexicographic maps $H_t:z\in \HH_+^*\mapsto F(z+\sqrt{-1}t)\in \HH_+^*$. It is clear that the induced action of $F$ on the  homology $H^{cl}_k(\conf_N(U))$ is that of $f$. 

The  action  of $f_o$ on $\Hcal_{N_0}(U_o)$ is through its action on  $\Ical_{U_o}|_{n_0}$ as a ring automorphism which  fixes the image of $\zeta$. Recall that this invokes (via Theorem \ref{thm:generation}) the K\"unneth products of  $f_o$-equivariant maps 
$\tri^\Ibold: \Ical_{U_o}|_i\to \Hcal_I(U_o)$ with $I\subset N_0$ and $i=|I|$.  
Hence  the same is true for its action  $\Cscr_{N}(U)_{n+k}$,  where we must take   $i=|I|\le n-k$. 

On the other hand,  we know that $f_o\in \Htp(U_o,*)$ acts on $H^{cl}_{n+k}(\conf_N(U))$ through its image 
$f\in \Htp(U,*)$. It is  geometrically  clear how: the maps $\tri^\Ibold: \Ical_{U_o}|_i\to \Hcal_I(U_o)$ involved in $\Cscr_{N}(U)_{n+k}$,  will after post composition  with the natural map  $\Hcal_I(U_o)\to \Hcal_I(U)$ factor through $\tri^\Ibold: \Ical_{U}|_i\to \Hcal_I(U)$. (More formally,  the image of  
of $\zeta-1 $ in  $\Ical_{U_o}|_i$ is mapped by $\tri^\Ibold$ to a  boundary in  the complex $\Cscr_{N}(U)_\pt$.)
This proves that $f$ acts on $H^{cl}_{n+k}(\conf_N(U))$ via its action on $\Ical_{U}|_{n-k}$.
\end{proof}

\end{document}
